% ===========================================================================
% KẾT LUẬN
% ===========================================================================
\chapter*{KẾT LUẬN}
\addcontentsline{toc}{chapter}{KẾT LUẬN}
\markboth{KẾT LUẬN}{KẾT LUẬN}

\section*{1. Kết quả nghiên cứu}
\addcontentsline{toc}{section}{1. Kết quả nghiên cứu}

Đề án đã hoàn thành mục tiêu nghiên cứu, thiết kế và xây dựng thử nghiệm mô hình
học máy hỗ trợ kiểm duyệt tín dụng. Các kết quả chính đạt được
gồm năm nhóm.

\subsection*{1.1. Về cơ sở lý luận}

Đề án đã hệ thống hóa cơ sở lý thuyết về tín dụng, kiểm duyệt tín dụng và đánh
giá rủi ro tín dụng, trong đó làm rõ hai nội dung có ý nghĩa định hướng cho toàn
bộ thiết kế. Thứ nhất là sự phân biệt giữa \textit{kiểm duyệt} -- quá trình đánh
giá, sàng lọc hồ sơ -- với \textit{phê duyệt} -- hành vi ra quyết định của cấp có
thẩm quyền; sự phân biệt này xác định vai trò hỗ trợ chứ không thay thế của hệ
thống. Thứ hai là việc xác định vị trí của bài toán trong khung lý thuyết rủi ro
tín dụng: trong phân rã $\EL = \PD \times \LGD \times \EAD$, đề án tập trung ước
lượng thành phần $\PD$.

Đề án cũng đã phân tích bốn hạn chế mang tính cấu trúc của phương pháp kiểm duyệt
truyền thống, ba đặc thù của bài toán tín dụng so với bài toán phân loại thông
thường, và trình bày nguyên lý cùng ưu nhược điểm của năm nhóm thuật toán học máy
được nghiên cứu.

\subsection*{1.2. Về quy trình xử lý dữ liệu}

Đề án đã xây dựng quy trình hoàn chỉnh từ thu thập dữ liệu tới tập đặc trưng sẵn
sàng cho mô hình, với ba đóng góp cụ thể:

\begin{itemize}
  \item \textbf{Bộ ràng buộc kiểm tra chất lượng dữ liệu} gồm \NConstraints{}
  ràng buộc nghiệp vụ, kèm cơ chế kiểm chứng bằng phép tiêm lỗi nhằm xác nhận
  bản thân cơ chế kiểm tra hoạt động đúng.

  \item \textbf{Thuật toán tính đặc trưng bảo đảm tính đúng đắn theo thời điểm}
  bằng kỹ thuật luỹ kế kết hợp ghép gần nhất theo thời gian, đưa chi phí tính
  toán từ $O(N \times |W| \times \bar{m})$ xuống tuyến tính. Quan trọng hơn chi
  phí, phương pháp này bảo đảm tính đúng đắn bằng chính cấu trúc thuật toán chứ
  không phụ thuộc vào kiểm tra bổ sung.

  \item \textbf{Quy trình lựa chọn đặc trưng ba bước} theo năng lực dự báo, mức
  tương quan và độ ổn định, rút gọn tập đặc trưng từ \NFeaturesInitial{} xuống
  \NFeaturesSelected{} biến.
\end{itemize}

\subsection*{1.3. Về mô hình học máy}

Đề án đã huấn luyện, hiệu chuẩn và so sánh \NModels{} mô hình trên cùng một điều
kiện thực nghiệm. Mô hình được chọn triển khai là \BestModel{}, đạt
$\AUC = \BestAUC{}$, $\text{Gini} = \BestGini{}$, $\KS = \BestKS{}$ và sai số
hiệu chuẩn kỳ vọng \BestECE{}.

Ngoài các con số cụ thể, thực nghiệm cho ba kết quả có giá trị nhận thức. Thứ
nhất, \textbf{khoảng cách giữa các thuật toán là nhỏ}: chênh lệch $\AUC$ giữa mô
hình được chọn và hồi quy logistic chỉ \GapBestLR{}, trong khi hồi quy logistic
có ưu thế rõ rệt về khả năng giải thích, chi phí huấn luyện và chất lượng hiệu
chuẩn tự nhiên. Thứ hai, \textbf{bảng xếp hạng $\AUC$ có giới hạn độ phân giải}:
phép đo dao động khoảng $3 \times 10^{-4}$ giữa các lần chạy, tức cùng bậc với
khoảng cách giữa các mô hình dẫn đầu, nên thứ hạng ở chữ số thập phân thứ tư
không phải bằng chứng về sự vượt trội. Đề án vì vậy chọn mô hình triển khai theo
tiêu chí biên Pareto -- nhanh nhất trong nhóm tương đương về chất lượng -- thay
vì lấy thẳng mô hình đứng đầu bảng. Thứ ba, \textbf{khoảng cách tới trần Bayes
cho biết giới hạn của việc cải tiến thuật toán}: mô hình được chọn còn cách trần
Bayes \BayesAUC{} khoảng
\GapBayes{} điểm $\AUC$, và phần chênh lệch này bắt nguồn từ việc dữ liệu quan sát được
chỉ là phép đo có nhiễu của mức tín nhiệm thật, chứ không phải do thuật toán yếu.
Hệ quả thực tiễn là cải thiện đáng kể chỉ đến từ việc bổ sung nguồn dữ liệu mới,
không phải từ thuật toán phức tạp hơn.

\subsection*{1.4. Về giá trị của các nguồn dữ liệu}

Phân tích SHAP cho thấy nhóm đặc trưng hành vi giao dịch đóng góp
\ShapBehaviourPct{}\% vào quyết định của mô hình, cao hơn nhóm lịch sử tín dụng
CIC (\ShapCICPct{}\%) vốn được coi là nguồn thông tin truyền thống quan trọng
nhất. Đây là căn cứ định lượng cho khuyến nghị khai thác dữ liệu giao dịch trong
công tác kiểm duyệt tín dụng: nguồn dữ liệu này do tổ chức tín dụng tự sở hữu,
không phát sinh chi phí truy vấn bên ngoài, phản ánh tình trạng tài chính hiện
tại và luôn có sẵn kể cả với nhóm khách hàng thiếu lịch sử tín dụng.

\subsection*{1.5. Về quy trình hỗ trợ kiểm duyệt}

Đề án đã hiện thực quy trình hoàn chỉnh kết hợp \NRules{} luật nghiệp vụ với mô
hình học máy theo logic ưu tiên tám mức, đóng gói thành dịch vụ REST với năm điểm
cuối hoạt động đúng. Mô phỏng trên tập kiểm thử cho tỷ lệ nợ xấu thực tế tăng đơn
điệu qua ba nhóm khuyến nghị (\BadApprove{} -- \BadReview{} -- \BadReject{}), tỷ
lệ xử lý tự động \AutomationRate{} và thời gian xử lý một hồ sơ ở phân vị 95\% là
\SysPNinetyFive{} ms.

So với quy trình thẩm định truyền thống chỉ dựa trên luật nghiệp vụ, mô hình
giảm tỷ lệ nợ xấu trong nhóm hồ sơ được duyệt từ \TradBadRate{} xuống
\PropBadRate{}, tức giảm tương đối \BadRateReductionPct{}. Đổi lại, tỷ lệ duyệt
giảm \ApproveRateDrop{} điểm phần trăm và khối lượng thẩm định thủ công tăng
\ReviewRateRise{} điểm phần trăm. Đánh đổi này điều chỉnh được thông qua hai
ngưỡng phân vùng, trong khi bộ quy tắc cố định không có cơ chế tương đương.

\section*{2. Hạn chế của nghiên cứu}
\addcontentsline{toc}{section}{2. Hạn chế của nghiên cứu}

Các hạn chế của kết quả thực nghiệm đã được phân tích chi tiết tại
mục~\ref{subsec:han-che-dinh-huong}. Ở tầm tổng thể, đề án có ba hạn chế cơ bản.

\textit{Thứ nhất là hạn chế về dữ liệu.} Do đặc thù bảo mật của dữ liệu tín dụng
cá nhân, đề án sử dụng dữ liệu mô phỏng sinh theo một mô hình xác suất do chính
đề án thiết kế, nên cấu trúc phụ thuộc giữa các biến phản ánh giả định của người
thiết kế chứ không phải quy luật thực tế. Giá trị tuyệt đối của các chỉ tiêu chất
lượng vì vậy không có ý nghĩa dự báo về hiệu quả trên dữ liệu thực; giá trị của
thực nghiệm nằm ở việc so sánh \textit{tương đối} giữa các mô hình trong cùng
điều kiện và ở việc kiểm chứng tính đúng đắn của toàn bộ quy trình.

\textit{Thứ hai là hạn chế về phương pháp.} Đề án chưa xử lý bài toán suy luận
nhãn cho hồ sơ bị từ chối (\eng{reject inference}) nên có thể chịu sai lệch chọn
mẫu; chưa thực hiện kiểm định công bằng thuật toán; và việc đánh giá độ ổn định
chỉ dựa trên bốn tháng dữ liệu kiểm thử, quá ngắn so với chu kỳ suy giảm điển
hình của mô hình tín dụng.

\textit{Thứ ba là hạn chế về phạm vi.} Đề án chỉ mô hình hóa thành phần $\PD$
trong phân rã tổn thất kỳ vọng; $\LGD$ và $\EAD$ không thuộc phạm vi nghiên cứu.
Mô hình được xây dựng ở mức nguyên mẫu trong môi trường nghiên cứu, chưa triển
khai trên hệ thống nghiệp vụ thực tế của một tổ chức tín dụng, và các số đo thời
gian chỉ có giá trị so sánh tương đối giữa các mô hình.

\section*{3. Hướng phát triển}
\addcontentsline{toc}{section}{3. Hướng phát triển}

Trên cơ sở các hạn chế đã nhận diện, đề án đề xuất năm hướng phát triển tiếp
theo, sắp xếp theo mức độ ưu tiên:

\begin{enumerate}[label=(\arabic*)]
  \item \textbf{Kiểm chứng trên dữ liệu thực} của một tổ chức tín dụng. Đây là
  hướng có giá trị cao nhất vì trực tiếp giải quyết hạn chế cơ bản nhất của đề
  án; mã nguồn đã được thiết kế để chuyển sang dữ liệu thực chỉ bằng cách thay
  thành phần tiếp nhận dữ liệu.

  \item \textbf{Bổ sung xử lý suy luận nhãn cho hồ sơ bị từ chối}, mở rộng phạm
  vi áp dụng của mô hình tới toàn bộ hồ sơ nộp vào.

  \item \textbf{Mở rộng sang mô hình $\LGD$ và $\EAD$} để hoàn thiện khung tính
  tổn thất kỳ vọng, cho phép định giá theo rủi ro đầy đủ.

  \item \textbf{Bổ sung kiểm định công bằng thuật toán và cơ chế cảnh báo tự
  động} khi phát hiện suy giảm chất lượng hoặc chênh lệch giữa các nhóm khách
  hàng.

  \item \textbf{Nghiên cứu cơ chế cập nhật mô hình có kiểm soát}, gồm tiêu chí
  kích hoạt huấn luyện lại, quy trình kiểm định mô hình mới và cơ chế quay lui
  khi phát hiện vấn đề trong vận hành.
\end{enumerate}

\vspace{0.8cm}

Tóm lại, đề án đã xây dựng được một mô hình học máy hỗ trợ kiểm duyệt tín dụng
hoàn chỉnh từ khâu tiếp nhận dữ liệu tới khâu đưa ra khuyến nghị kèm lý giải,
đồng thời đánh giá định lượng giá trị mà mô hình mang lại so với quy trình thẩm
định
truyền thống. Kết quả nghiên cứu cho thấy việc kết hợp mô hình học máy với bộ
luật nghiệp vụ, trong đó mô hình đóng vai trò cung cấp thông tin định lượng còn
quyết định cuối cùng vẫn thuộc về người có thẩm quyền, là hướng tiếp cận vừa khả
thi về kỹ thuật vừa phù hợp với yêu cầu quản trị rủi ro và tuân thủ của hoạt động
cấp tín dụng.
